\documentclass[12pt]{article}
\usepackage[utf8]{inputenc}
\usepackage[T1]{fontenc}
\usepackage[margin=1in]{geometry}
\usepackage{setspace}
\usepackage{xcolor}
\usepackage[colorlinks=true, urlcolor=blue!60!black, linkcolor=blue!60!black]{hyperref}
\setlength{\parindent}{0em}
\setlength{\parskip}{0.6em}
\renewcommand{\baselinestretch}{1.0}
\begin{document}
\begin{center}
{\LARGE \textit{City of Boerne v. Flores}, 521 U.S. 507 (1997)\par}
\vspace{0.5em}
{\large Case Brief --- Two-Page Edition\par}
\end{center}

\section{Memory Jogger}
Striking down RFRA as applied to the states, the Court held that Section 5 of the Fourteenth Amendment is a remedial power requiring ``congruence and proportionality,'' not a power to redefine rights (\href{https://supreme.justia.com/cases/federal/us/521/507/}{521 U.S. 507}).

\section{Detailed Case Facts}
Boerne, Texas, denied Archbishop Flores a permit to enlarge St. Peter Catholic Church because it lay in a new historic district. Flores sued under RFRA, which Congress enacted after \textit{Employment Division v. Smith} (1990) to ``restore the compelling interest test'' against all governments, relying on Section 5 to bind the states.

\section{Procedural History}
The district court held RFRA beyond Section 5; the Fifth Circuit reversed; the Supreme Court granted certiorari and reversed.

\section{Judicial Votes}
Reversed, 6--3. Kennedy wrote for Rehnquist, Stevens, Thomas, Ginsburg, and (except Part III-A-1) Scalia; Stevens and Scalia concurred; O'Connor, Souter, and Breyer dissented.

\section{Holding}
As applied to state and local governments, RFRA exceeds Congress's Section 5 power. Congress may enforce, but not alter, constitutional rights, and remedial legislation must be congruent and proportional to a documented injury.

\section{Analysis of Opinions}
Kennedy held that Section 5 grants only the power ``to enforce''; the 39th Congress rejected a draft granting plenary power, and letting Congress redefine rights would circumvent Article V. Unlike the Voting Rights Act, RFRA rested on no record of modern religious bigotry, lacked sunset or geographic limits, and imposed ``the most demanding test known to constitutional law.'' Stevens found an Establishment Clause problem; Scalia defended \textit{Smith}'s history. O'Connor accepted the Section 5 framework but would have overruled \textit{Smith}; Souter and Breyer would have ordered reargument.

\section{Examples: Future Applications}
\textbf{Same-Side:} ADA damages suits against states for failing to accommodate disabilities, unsupported by a record of unconstitutional discrimination, fail (\textit{Garrett}). \textbf{Opposite-Side:} a narrow family-leave mandate answering documented gender discrimination is upheld (\textit{Hibbs}). \textbf{Fence-Sitter:} a nationwide ban on voter-ID rules with racial disparities, without sunset or geographic limits.

\section{Critique}
Post and Siegel attack the decision's ``juricentric'' model, and McConnell argues that ``appropriate legislation'' and Reconstruction history support deference to Congress. Federalism defenders such as Jeffrey Sutton reply that without \textit{Boerne} Congress would gain a general police power and could amend the Constitution without Article V.

\section{Key Quotations}
\begin{quote}
``There must be a congruence and proportionality between the injury to be prevented or remedied and the means adopted to that end. Lacking such a connection, legislation may become substantive in operation and effect.''
\end{quote}
\noindent --- Justice Kennedy, Opinion of the Court, 521 U.S. at 520.

\end{document}
